\section{Основы машинного обучения и искусственных нейронных сетей}
\label{sec:ml}

\noindent\hspace{1.25cm}\emph{Раздел вводит понятия, общие для ЛР~№1--4.}

\subsection{Задачи машинного обучения}

\emph{Машинное обучение}~--- построение алгоритмов, которые извлекают закономерность из данных, а не
получают её в виде явно запрограммированных правил. Пусть имеется \emph{обучающая выборка}
$\{(\mathbf{x}_k, \mathbf{t}_k)\}_{k=1}^{N}$, где $\mathbf{x}_k \in \mathbb{R}^n$~--- вектор признаков
(вход), $\mathbf{t}_k$~--- известный ответ (цель, target). Требуется построить модель $\mathbf{y} =
F(\mathbf{x}; \boldsymbol\theta)$ с параметрами $\boldsymbol\theta$, которая даёт правильные ответы не
только на обучающих, но и на новых примерах. Способность модели работать на новых данных называется
\emph{обобщением}~\cite{golovko2017, haykin2006}.

В лабораторном практикуме встречаются четыре типа задач:
\begin{itemize}
\item \emph{классификация}~--- ответ принадлежит конечному множеству классов (ЛР~№1): распознать, к
какой из групп относится точка плоскости;
\item \emph{регрессия (аппроксимация)}~--- ответ вещественный (ЛР~№3): восстановить функцию $y = f(t)$ по
её значениям в узлах сетки;
\item \emph{прогнозирование временного ряда}~--- частный случай регрессии, где входом служат прошлые
значения сигнала, а целью~--- следующее значение (ЛР~№2);
\item \emph{кластеризация}~--- разбиение объектов на группы похожих без известных ответов (ЛР~№4).
\end{itemize}

Первые три задачи решаются \emph{обучением с учителем} (supervised learning): для каждого входа известен
правильный выход, и параметры подстраиваются так, чтобы уменьшить расхождение между ответом модели и
целью. Кластеризация относится к \emph{обучению без учителя} (unsupervised learning): модель извлекает
структуру только из взаимного расположения входов~\cite{ablameyko2011}. Генетические алгоритмы (разделы~5--7) формально не
являются методом обучения, но решают ту же вычислительную задачу~--- поиск параметров, доставляющих
экстремум целевой функции.

\subsection{Функция потерь, метрики и обобщение}

Качество модели на выборке измеряется \emph{функцией потерь}. Для регрессии стандартной является
среднеквадратичная ошибка (mean squared error, MSE)
\begin{equation}
\label{eq:mse}
E(\boldsymbol\theta) = \frac{1}{N}\sum_{k=1}^{N} \bigl(y_k - t_k\bigr)^2 ,
\end{equation}
в Python~--- \code{np.mean((y - t)**2)} или функция \code{mean\_squared\_error} из \code{sklearn.metrics}. Для
интерпретации результата используют коэффициент детерминации $R^2 = 1 - E / \sigma_t^2$
(\code{r2\_score}), где
$\sigma_t^2$~--- дисперсия целевых значений: $R^2 = 1$ означает точное совпадение, $R^2 = 0$~---
модель не лучше константы. Для классификации применяют долю верных ответов (accuracy) и число
ошибок на эпоху.

Обучение~--- это минимизация $E(\boldsymbol\theta)$, но цель~--- малая ошибка на \emph{новых} данных.
Поэтому выборку делят на три части: \emph{обучающую} (по ней настраиваются параметры),
\emph{валидационную} (по ней выбирают архитектуру и момент остановки) и \emph{тестовую} (для
окончательной независимой оценки), обычно в пропорции 70/15/15\,\%; в Python случайное разбиение
выполняет \code{sklearn.model\_selection.train\_test\_split}.

Соотношение ошибок на обучении и на новых данных определяется \emph{ёмкостью} модели~--- числом и
гибкостью её параметров:
\begin{itemize}
\item \emph{недообучение} (underfitting)~--- модель слишком проста, ошибка велика и на обучающих, и на
новых данных, форма зависимости не передаётся;
\item \emph{переобучение} (overfitting)~--- модель слишком гибка и подстраивается под случайный шум
обучающей выборки: ошибка на обучении мала, на новых данных велика.
\end{itemize}
Ожидаемая ошибка на новых данных раскладывается на квадрат смещения (систематическая ошибка простой
модели), дисперсию (чувствительность сложной модели к конкретной выборке) и неустранимый шум данных.
С ростом ёмкости смещение убывает, а дисперсия растёт, поэтому зависимость ошибки на новых данных от
сложности модели имеет U-образную форму (подраздел~\ref{sec:generalization}). Способы борьбы с переобучением~--- уменьшение числа параметров,
регуляризация (штраф за большие веса), увеличение выборки и \emph{ранний останов}: обучение
прекращают, когда ошибка на валидационной выборке перестаёт уменьшаться.

\subsection{Модель искусственного нейрона}

Искусственный нейрон (рисунок~\ref{fig:neuron})~--- упрощённая модель нервной клетки, предложенная
У.~Мак-Каллоком и У.~Питтсом в 1943~г. Нейрон получает входы $x_1, \ldots, x_n$, умножает их на
\emph{синаптические веса} $w_1, \ldots, w_n$, суммирует вместе со \emph{смещением} (bias) $b$ и
пропускает результат через \emph{функцию активации}~$f$:
\begin{equation}
\label{eq:neuron}
z = \sum_{i=1}^{n} w_i x_i + b = \mathbf{w}^{\mathsf T}\mathbf{x} + b, \qquad y = f(z).
\end{equation}
Величина $z$ называется \emph{сетевым входом} (net input). Смещение удобно считать весом
дополнительного входа, постоянно равного единице.

\begin{figure}[!htb]
\centering
\begin{tikzpicture}[>=Stealth, node distance=8mm,
  inp/.style={circle, draw, minimum size=7mm, fill=cblue!12},
  op/.style={circle, draw, minimum size=10mm, thick}]
\node[inp] (x1) at (0, 1.6) {$x_1$};
\node[inp] (x2) at (0, 0.55) {$x_2$};
\node at (0, -0.35) {$\vdots$};
\node[inp] (xn) at (0, -1.3) {$x_n$};
\node[inp, fill=clight] (one) at (1.6, -2.3) {$1$};
\node[op] (sum) at (4, 0) {$\Sigma$};
\node[op, rounded rectangle, minimum width=14mm] (f) at (6.4, 0) {$f(z)$};
\draw[->] (x1) -- node[above, sloped, font=\small] {$w_1$} (sum);
\draw[->] (x2) -- node[above, sloped, font=\small] {$w_2$} (sum);
\draw[->] (xn) -- node[below, sloped, font=\small] {$w_n$} (sum);
\draw[->] (one) -- node[below right, font=\small] {$b$} (sum);
\draw[->] (sum) -- node[above, font=\small] {$z$} (f);
\draw[->] (f) -- ++(1.8, 0) node[right] {$y$};
\node[font=\small, align=center, text=cgray] at (4, 1.5) {сумматор};
\node[font=\small, align=center, text=cgray] at (6.4, 1.1) {функция\\активации};
\end{tikzpicture}
\caption{Модель искусственного нейрона}
\label{fig:neuron}
\end{figure}

Вид функции активации определяет свойства нейрона (рисунок~\ref{fig:activation}):
\begin{itemize}
\item пороговая (\code{hardlim}): $f(z) = 1$ при $z \ge 0$ и $f(z) = 0$ при $z < 0$, в NumPy
\code{(z >= 0).astype(int)}~--- персептрон Розенблатта (раздел~\ref{sec:perceptron});
\item линейная (\code{purelin}): $f(z) = z$~--- линейная сеть ADALINE и выходной слой сети для регрессии;
\item логистическая сигмоида (\code{logsig}): $f(z) = 1/(1 + e^{-z})$, значения в $(0; 1)$,
в SciPy~--- \code{scipy.special.expit};
\item гиперболический тангенс (\code{tansig}): $f(z) = \tanh z = 2/(1 + e^{-2z}) - 1$, значения в
$(-1; 1)$, \code{np.tanh}~--- скрытый слой многослойного персептрона (раздел~\ref{sec:mlp}).
\end{itemize}
Сигмоидальные функции гладкие, монотонные и ограниченные, а их производные выражаются через значение
самой функции: $\sigma'(z) = \sigma(z)(1 - \sigma(z))$ для \code{logsig} и $\tanh'(z) = 1 - \tanh^2 z$.
Это делает вычисление градиентов при обучении дешёвым. В глубоких сетях сейчас чаще применяют
кусочно-линейную функцию ReLU $f(z) = \max(0, z)$, но для небольших сетей практикума классические
сигмоиды остаются стандартом.

\begin{figure}[!htb]
\centering
\includegraphics[width=\textwidth]{\figdir/activation.png}
\caption{Функции активации нейрона}
\label{fig:activation}
\end{figure}

\subsection{Архитектуры нейронных сетей}

Нейроны объединяются в \emph{слои}: все нейроны слоя получают один и тот же входной вектор и работают
параллельно. Слой из $S$ нейронов с $n$ входами описывается матрицей весов $W$ размера $S \times n$ и
вектором смещений $\mathbf{b}$ длины $S$: $\mathbf{y} = f(W\mathbf{x} + \mathbf{b})$. По способу
соединения слоёв различают:
\begin{itemize}
\item \emph{сети прямого распространения} (feedforward): сигнал проходит от входа к выходу без обратных
связей; к ним относятся персептрон, линейная сеть и многослойный персептрон (ЛР~№1--3);
\item \emph{рекуррентные сети}: выходы подаются обратно на вход с задержкой, что позволяет
обрабатывать последовательности;
\item \emph{сети с конкурентным слоем}: нейроны слоя соревнуются за право ответить на вход, как в картах
Кохонена (ЛР~№4).
\end{itemize}

В таблице~\ref{tab:nn-overview} сведены модели, изучаемые в практикуме.

\begin{table}[!htb]
\caption{Нейросетевые модели лабораторного практикума}
\label{tab:nn-overview}
\centering\small
\begin{tabularx}{\textwidth}{>{\raggedright}p{2.9cm}>{\raggedright}p{3.2cm}>{\raggedright}p{3.4cm}>{\raggedright\arraybackslash}X}
\toprule
Модель & Обучение & Задача & Библиотечный аналог для сравнения \\
\midrule
Персептрон & правило персептрона, с учителем & классификация линейно разделимых данных & \code{sklearn.linear\_model.} \code{Perceptron} \\
Линейная сеть & МНК, правило Уидроу~--- Хоффа & прогнозирование, линейная регрессия & \code{np.linalg.lstsq}, \code{LinearRegression}, \code{SGDRegressor} \\
Многослойный персептрон & обратное распространение ошибки & аппроксимация, классификация & \code{MLPRegressor}, \code{torch.nn} \\
Карта Кохонена & конкурентное, без учителя & кластеризация, визуализация & \code{MiniSom}, \code{KMeans} \\
\bottomrule
\end{tabularx}
\end{table}

Обучение сети выполняется \emph{эпохами}: за одну эпоху сети предъявляются все примеры обучающей
выборки. При \emph{последовательном} (стохастическом, онлайн) обучении веса меняются после каждого
примера, при \emph{пакетном}~--- один раз за эпоху по суммарному градиенту; промежуточный вариант~---
обучение \emph{мини-пакетами}, стандартное для PyTorch (\code{DataLoader(batch\_size=\ldots)}).

\subsection{Вычислительные инструменты Python}
\label{sec:python}

Все алгоритмы пособия реализуются на Python с библиотекой NumPy, а готовые библиотеки используются для
проверки и сравнения собственной реализации~\cite{harris2020, virtanen2020, pedregosa2011}. Роли
основных библиотек приведены в таблице~\ref{tab:python}.

\begin{table}[!htb]
\caption{Библиотеки Python, используемые в практикуме}
\label{tab:python}
\centering\small
\begin{tabularx}{\textwidth}{>{\raggedright}p{3.4cm}>{\raggedright\arraybackslash}X}
\toprule
Библиотека & Назначение \\
\midrule
NumPy & массивы, векторные операции, линейная алгебра (\code{np.linalg}), генератор случайных чисел \\
SciPy & эталонные решения: \code{optimize.minimize\_scalar}, \code{differential\_evolution}, \code{least\_squares} \\
scikit-learn & готовые модели и метрики: \code{Perceptron}, \code{MLPRegressor}, \code{KMeans}, \code{r2\_score} \\
PyTorch & нейронные сети с автоматическим дифференцированием (\code{torch.nn}, \code{torch.optim}) \\
MiniSom & самоорганизующиеся карты Кохонена \\
DEAP, PyGAD & библиотеки эволюционных вычислений (только для сравнения) \\
Matplotlib, Plotly, ipywidgets & графики, анимация, ползунки и кнопки в Jupyter \\
\bottomrule
\end{tabularx}
\end{table}

Три соглашения делают код быстрым и воспроизводимым. \emph{Векторизация}: выборка или популяция
хранится одним массивом (строка~--- объект, столбец~--- признак или ген), и операции выполняются над
массивом целиком без циклов Python. \emph{Транслирование} (broadcasting): массивы разной формы
согласуются автоматически, например \code{X[:, None] - X[None]} даёт все попарные разности.
\emph{Воспроизводимость}: все случайные числа берутся из одного генератора
\code{rng = np.random.default\_rng(seed)}, и значение \code{seed} сохраняется вместе с результатом.
